\documentclass[12pt]{article}

% Packages
\usepackage[utf8]{inputenc}
\usepackage[T1]{fontenc}
\usepackage{mathptmx} % Times New Roman font
\usepackage{geometry}
\usepackage{setspace}
\usepackage{titlesec}

% Page geometry
\geometry{a4paper, margin=1in}

% Single spacing
\singlespacing

% Section formatting - remove numbering and adjust spacing
\titleformat{\section}{\normalfont\bfseries}{\thesection}{0em}{}
\titlespacing*{\section}{0pt}{12pt}{6pt}

\begin{document}

% Title
\begin{center}
    \Large\textbf{Low Power Edge Inference for Green Computing: An Evaluation of Post-Training Quantization and Unstructured Pruning}
\end{center}

\vspace{12pt}

% Author Information - Two Column Layout
\begin{center}
\begin{tabular}{cc}
\textbf{Harshpreet Singh} & \textbf{Sameer Kumar}\\[3pt]
Chitkara University, India & Chitkara University, India\\[3pt]
\texttt{harshpreet393.be22@chitkara.edu.in} & \texttt{sameer776.be22@chitkara.edu.in}
\end{tabular}
\end{center}

\vspace{18pt}

% Abstract Title
\begin{center}
    \textbf{Abstract}
\end{center}

\vspace{6pt}

% Structured Abstract Content
\noindent\textbf{Background/Introduction}\\
With data center AI workloads projected to exceed 945 TWh annually by 2030, edge computing offers a sustainable alternative for deploying language models locally. However, practitioners lack clear guidance on which post-training quantization (PTQ) method best balances accuracy and energy efficiency for specific models like Qwen3-4B and Phi-2.

\vspace{12pt}

\noindent\textbf{Objectives}\\
This study compares four PTQ techniques, INT8 uniform quantization, GPTQ, AWQ, and INT4 weight-only on Qwen3-4B (4B) and Phi-2 (2.7B) to identify optimal accuracy-energy trade-offs for edge deployments.

\vspace{12pt}

\noindent\textbf{Methods/Methodology}\\
We evaluated all four PTQ approaches at 4-bit to 8-bit precision on Mac M4 Pro with 24GB unified memory. We measured MMLU accuracy, latency, memory footprint, and energy consumption per inference on both models.

\vspace{12pt}

\noindent\textbf{Results/Findings}\\
Initial observations suggest AWQ maintains better accuracy for Qwen3-4B at 4-bit precision, while INT8 proves reliable across both models. GPTQ appears more effective for Phi-2 than Qwen3-4B, and naive INT4 shows significant degradation. Model size influences quantization tolerance, with smaller architectures handling aggressive compression better. AWQ shows promise for Qwen3-4B and INT8 for Phi-2, though comprehensive benchmarking is ongoing.

\vspace{12pt}

\noindent\textbf{Conclusion}\\
Optimal PTQ strategy depends on model size: AWQ benefits larger models like Qwen3-4B through activation-awareness, while Phi-2 achieves excellent efficiency with simpler INT8 quantization, enabling tailored green AI deployment.

\vspace{18pt}

% Keywords
\noindent\textbf{Keywords:} Post-training quantization, unstructured pruning, edge AI, green computing, model compression, sustainable machine learning

\end{document}
